\subsection{A coefficient and discriminant form of the remaining level-$27$ test}

\begin{proposition}[B-RD36: the tame character is read from the discriminant]
\label{prop:B-RD36}
\status{Proved. Complete root and independent proof, original-source, and exact-certificate audits passed.}
Let $K$ have algebraically closed residue field of characteristic three,
and let $E/K$ be an elliptic curve with potentially supersingular
reduction. For any Weierstrass equation put
$\delta=v_K(\Delta)$. The order $t$ of its tame special-inertia
quotient is
\[
                  t=\frac4{\gcd(4,\delta)}.
                                                        \tag{B-RD36.1}
\]
Here $\gcd(4,0)=4$, and changing the equation changes $\delta$ by
a multiple of twelve. Hence the three possibilities are
\[
 \begin{array}{c|ccc}
 \delta\pmod4&0&2&1\text{ or }3\\ \hline
 t&1&2&4.
 \end{array}
                                                        \tag{B-RD36.2}
\]
Together with (B-RD32.2), this determines the requested
prime-to-$3$ character as a function on inertia.

For absolute ramification $e_K=2$, the unresolved alternative
$\nu_{\min}(3,2,2)=3$ is consequently equivalent to the existence,
for some algebraically closed $k$ of characteristic three, of two
triples $(a_2,a_4,a_6)$ in
\[
 R_k=W(k)[\pi]/(\pi^2-3),\qquad K_k=\operatorname{Frac}(R_k),
                                                        \tag{B-RD36.3}
\]
such that the integral equations
$y^2=x^3+a_2x^2+a_4x+a_6$ have nonzero discriminants, satisfy
\[
        0\leq v(\Delta)\leq23,\qquad
        3v(c_4)>v(\Delta),
                                                        \tag{B-RD36.4}
\]
have isomorphic generic $27$-torsion group schemes, and their
discriminant valuations belong to different columns of
(B-RD36.2). The convention is $v(0)=+\infty$.
This is a six-coefficient elliptic test. It is not asserted here
to be a finite residue-field elimination that has been evaluated.
\end{proposition}
\begin{proof}
Take a minimal faithful Galois good-reduction extension $L/K$.
Its group has order $m=3^wt$, with $w\in\{0,1\}$ and
$t\in\{1,2,4\}$, by the norm-one classification and the
canonical alternating pairing of an elliptic curve. A change
from the given Weierstrass equation to a good equation over
$\mathcal O_L$ has scaling factor $u$, with
\[
                  r=v_L(u)=m\delta/12\in\mathbf Z.
                                                        \tag{B-RD36.5}
\]
Choose a tame element $\tau$ of order $t$. Its action on a
uniformizer has residual multiplier a primitive $t$-th root
$\xi$. The residual multiplier of $\tau(u)/u$ is therefore
$\xi^r$. Up to inversion, it is the multiplier on the invariant
differential of the good special elliptic curve. By the explicit
automorphisms (B-RD35.1), the tangent character has precisely the
wild order-three group as its kernel. Thus the multiplier of
$\tau$ has order exactly $t$, and $\gcd(r,t)=1$.

If $t=1$, equation (B-RD36.5) gives $4\mid\delta$.
If $t>1$, $r$ is odd; taking two-adic valuations in
$12r=3^wt\delta$ gives
$2=v_2(t)+v_2(\delta)$. This proves (B-RD36.1)--(B-RD36.2).
The character identification follows from
Lemma~\ref{lem:B-RD32-norm}.

Every complete mixed-characteristic DVR of absolute ramification
two with algebraically closed residue field $k$ contains its Witt
coefficient ring. Write $3=\varpi^2v$ for a uniformizer $\varpi$
and a unit $v$. Since $k$ is algebraically closed and $2$ is
invertible, Hensel lifting gives a square root of $v$. Replacing
$\varpi$ by $\varpi\sqrt v$ gives (B-RD36.3); the degree-two
coefficient-ring description follows from the Cohen structure
argument used in Proposition~\ref{prop:B-RD24-finite-test}.
Such a field contains $\mu_3$, since its coefficient ring contains
a square root of $-1$.

By Proposition~\ref{prop:B-RD35}, every remaining BT comparison is
an elliptic comparison over one of these fields. Choose a minimal
integral Weierstrass equation and complete the square, which is an
integral change since $2$ is a unit. The discriminant bound in
Proposition~\ref{prop:B-RD24-finite-test}, with $p=3,e=2$,
gives $v(\Delta)\leq6e+11=23$.

We justify the strict $j$-condition in (B-RD36.4). Potential good
reduction is equivalent to integral $j$
\cite[Section~2, p.~496]{SerreTate1968Good}. On a good equation
in characteristic three the Cartier operator on
$\omega=dx/y$, for $y^2=x^3+a_2x^2+a_4x+a_6$, is
\[
                 C(\omega)=a_2^{1/3}\omega.
\]
Indeed $\omega=(x^3+a_2x^2+a_4x+a_6)\,dx/y^3$, and Cartier
selects the $x^2dx$ coefficient. Also $c_4\equiv a_2^2\pmod3$.
Thus a good elliptic special fibre is supersingular exactly when
its $j$-invariant is zero. Since $j=c_4^3/\Delta$, potential
supersingularity is exactly $3v(c_4)>v(\Delta)$, including $j=0$.

Conversely any pair in the stated coefficient test is potentially
supersingular and gives two allowed height-two BT groups. Its
different discriminant columns give different prime-to-$3$
characters, while its isomorphic $27$-torsion gives a failure at
manuscript index two. The proved upper bound in
Theorem~\ref{thm:B-RD34} then makes the least index three.
If no such pair exists, B-RD35 and B-RD34 instead make the least
index two. No existence or nonexistence of such a pair is inferred
from this reformulation.
\end{proof}

\paragraph{Exact finite certificate.}
The retained program
\path{code/verify_discriminant_tame_character.py} checks all twenty-four
displayed discriminant valuations against the necessary scaling equation
$12r=m\delta$ and the tame coprimality condition. It also checks 4,221
changes $\delta\mapsto\delta+12z$, 977 instances of the preceding crude
minimal-discriminant bound, and the exact polynomial identities
$c_4^3-c_6^2=1728\Delta$, $c_4\bmod3=a_2^2$, and the Cartier coefficient
$a_2$. The author run passed in $0.04732804800005397$ seconds;
\path{code/verify_discriminant_tame_character.output.json} retains its
full output. These finite checks certify the displayed arithmetic and
identities. They assert neither existence of all necessary scaling cases
nor existence of a level-$27$ comparison pair. Independent reruns by the
root reviewer and the second reviewer passed in respectively
$0.04524370399667532$ and $0.061048945997754345$ seconds; the author output
was preserved.

\paragraph{Subsequent exact determination.}
The remaining small-height alternatives isolated here are settled in
Theorem~\ref{thm:B-RD42}: the exact manuscript index is
$\nu_{\min}(3,2,2)=2$. The intermediate bounds and conditional
reductions above are retained as steps in that proof.
